\documentclass{article}
\usepackage[UTF8]{ctex}
\usepackage{amsmath,amssymb,amsfonts}
\usepackage{enumitem}
\usepackage{geometry}
\geometry{a4paper, margin=1in}
\usepackage{esint}
% === 错题本排版优化（不参与题目标识） ===
\setlength{\parindent}{0pt}
\setlength{\parskip}{0.35em plus 0.10em minus 0.08em}
\setlength{\emergencystretch}{3em}
\setlist{itemsep=0.25em, topsep=0.40em, parsep=0pt}
\allowdisplaybreaks[2]
% === 错题本排版优化结束 ===
\begin{document}
\section*{南京理工大学2026年数学分析真题}
\section*{97．南京理工大学2026年数学分析试卷（回忆版）}

\subsection{题目1}
1．解答如下问题：

\begin{itemize}
  \item[（1）] 利用定义证明：\((\cos x)^{\prime}=-\sin x\) ．
  \item[（2）] 利用定积分定义证明： \(\int_{0}^{1} x^{3} \mathrm{~d} x=\frac{1}{4}\) ．
\end{itemize}

\subsection{题目2}
2 ．求极限：

\begin{itemize}
  \item[（1）] \(\lim _{n \rightarrow \infty} \sum_{k=1}^{n} \frac{2^{\frac{k}{n}}}{n+\frac{1}{k}}\) ．
  \item[（2）] \(\lim _{x \rightarrow 0}\left(\frac{1}{\sin x}-\frac{\sin x}{e^{x}-1}\right)\) ．
\end{itemize}

\subsection{题目3}
3．设 \(f(x, y)=\left\{\begin{array}{ll}\frac{x y^{2}}{x^{2}+y^{4}}, & (x, y) \neq(0,0), \\ 0, & (x, y)=(0,0) .\end{array}\right.\) 讨论 \(f(x, y)\) 在点 \((0,0)\) 的连续性，偏导数的存在性，可微性。

\subsection{题目4}
4．求幂级数 \(\sum_{n=0}^{\infty} \frac{(n+1)(n+2)}{n+3} x^{n}\) 的收敛域以及和函数．

\subsection{题目5}
5．设 \(f(x)\) 在 \(I\) 上有连续的导数，记 \(F_{n}(x)=n f\left(x+\frac{1}{n}\right)-n f(x)\) ．

\begin{itemize}
  \item[（1）] 若 \(I\) 为有界闭区间，\(\left\{F_{n}(x)\right\}\) 是否一致收敛？给出理由．
  \item[（2）] 若 \(I\) 为有界开区间，\(\left\{F_{n}(x)\right\}\) 是否一致收敛？给出理由．
\end{itemize}

\subsection{题目6}
6．设 \(f(x)\) 为连续函数，求曲线积分

\[
\int_{L} \frac{1+y^{2} f(x y)}{y} \mathrm{~d} x+\frac{x\left(y^{2} f(x y)-1\right)}{y^{2}} \mathrm{~d} y .
\]

其中 \(L\) 为 \(\left(3, \frac{2}{3}\right)\) 指向 \((1,2)\) 的直线．

\subsection{题目7}
7．证明： \(\int_{0}^{+\infty} \frac{\sin x}{x} \mathrm{~d} x=\frac{\pi}{2}\) ，并求 \(\int_{0}^{+\infty} \frac{\sin ^{2}(x y)}{x^{2}} \mathrm{~d} x\) ．

\subsection{题目8}
8．求曲面积分

\[
\oiint_{\Omega} \frac{x}{\left(x^{2}+y^{2}+z^{2}\right)^{\frac{3}{2}}} \mathrm{~d} y \mathrm{~d} z+\frac{y}{\left(x^{2}+y^{2}+z^{2}\right)^{\frac{3}{2}}} \mathrm{~d} z \mathrm{~d} x+\frac{z}{\left(x^{2}+y^{2}+z^{2}\right)^{\frac{3}{2}}} \mathrm{~d} x \mathrm{~d} y .
\]

其中 \(\Omega\) 是区域 \(|x| \leq 3,|y| \leq 5,|z| \leq 4\) 的表面，方向向外。

\subsection{题目9}
9．解答如下问题：

\begin{itemize}
  \item[（1）] 叙述有限覆盖定理以及一致连续的定义．
  \item[（2）] 利用有限覆盖定理证明连续函数在闭区间内一致连续．
\end{itemize}
\end{document}
